\section{The two additive passage models}\label{sec:additive}

The two scalar passages below have an action given by a linear function
of the passage variable plus a logarithmic correction.  One describes
a damped transverse state; the other describes a residual state driven by layers near the two
endpoints.  The common small quantity is the exponential of the negative
action difference.  We construct packet expansions for both passages,
including the independent derivatives and retestability required by
Sections~\ref{sec:calculus} and~\ref{sec:elimination}.

The coefficients and ordinary parameters are analytic on neighborhoods
slightly larger than the boxes used below.  Amplitudes are chosen smaller
than a fixed constant determined by these boxes, independently of any
asymptotic order.  Constants for a fixed coefficient, charge budget, or
derivative order may depend on that request.

We first construct the actual analytic passage.  Its finite expansion in
the transfer charge will have coefficients built from slow sectorial
solutions and their primitives.  We retain their exact inner endpoint
values, then expand each charge coefficient successively at the flag
indices where its small arguments first occur.  These are three distinct
choices: a finite charge cutoff, an optimal slow Gevrey cutoff, and finite
Taylor degrees for the later flag transitions.  No convergence of the
entire packet tree is required.

\subsection{Models, independent inputs, and actual continuation}

The real inputs satisfy
\begin{equation}\label{eq:additive-inputs}
 q\ge q_*>1,\qquad Q/q\longrightarrow\infty.
\end{equation}
Here $q_*$ is chosen once from the field and its boxes.  Either $q$ tends
to infinity or it belongs to an ordinary neighborhood of a fixed positive
value.  Put
\begin{equation}\label{eq:additive-clock}
 \gamma=\delta q,\qquad A(w)=w+\gamma\log(w/q),\qquad
 A(p)=Q,\qquad P=e^{-(Q-q)},\qquad Z=\gamma/w,
\end{equation}
where $\delta$ is small and all logarithms have the stated lifted
determination.  Thus $p$ is the outer endpoint in the $w$ coordinate,
whereas $Q$ is its action value.  Since $A(q)=q$, the action difference
is $Q-q$ and its decaying exponential is $P$.  The two endpoints are
related by
\begin{equation}\label{eq:additive-endpoint-unit}
 p=QU,\qquad U+\delta(q/Q)\{\log(Q/q)+\log U\}=1.
\end{equation}
The solution $U$ near $1$ is analytic in the small ratio and
ratio--logarithm arguments appearing here.  We use a finite vector $X$ of
profiles
\begin{equation}\label{eq:additive-profiles}
 X_j=\chi_j(w/q)^{r_j}\quad(r_j<0),\qquad
 X_j=\chi_j(w/Q)^{r_j}\quad(r_j>0),
\end{equation}
with fixed rational exponents and small amplitudes.  The profile $Z$ may
be included among the negative profiles.  Redundant amplitudes and all
ordinary field parameters may first be taken as independent arguments.

The two models are
\begin{align}
 \text{(A)}\quad &v'=\sigma(1+Z)v+w^{-1}g(X,v),
       &&\sigma\in\{-1,1\},\quad g(0,v)=0,\label{eq:additive-A}\\
 \text{(B)}\quad &z'=G(X,E,H,z),
       &&G(X,0,0,z)=0,\label{eq:additive-B}
\end{align}
where prime denotes $\dd/\dd w$.  Model (A) is traversed from $q$ to $p$
when $\sigma=-1$ and from $p$ to $q$ when $\sigma=1$; its input lies in a
small disk about $0$.  Since $A'(w)=1+Z$, the homogeneous multiplier
in either prescribed direction is exactly $P$.  Model (B) is traversed
from $q$ to $p$ and has its input in a smaller disk about any chosen
ordinary reference state.  Its fast profiles are
\begin{equation}\label{eq:additive-fast}
 E_s=e_s(w/q)^{\beta_s}e^{-a_s(A(w)-q)},\qquad
 H_s=d_s(w/Q)^{\nu_s}e^{-b_s(Q-A(w))}.
\end{equation}
The exponents $\beta_s,\nu_s$ are rational, the amplitudes are small, and
the positive weights $a_s,b_s$ belong to one fixed discrete charge lattice.
We measure charge in its positive unit and write the weights as positive
integers; with another unit every occurrence of $P$ is replaced by the
corresponding fixed power of $P$.

Write $\eta$ for the ordinary field parameters suppressed in the equations.
The output maps have the independent arguments
\begin{align*}
 \text{(A)}\quad &(q,Q,\delta,\chi,v_{\rm in},\eta)
          \longmapsto v_{\rm out},\\
 \text{(B)}\quad &(q,Q,\delta,\chi,e,d,z_{\rm in},\eta)
          \longmapsto z_{\rm out}.
\end{align*}
The endpoint $p$ is determined by $A(p)=Q$.  Relations between amplitudes,
clocks, and physical endpoints are imposed only in the applications in
Section~\ref{sec:terminal}; they are not imposed when differentiating
these primitive families.  A flag test expresses the independent passage
arguments in free flag coordinates and ordinary coordinates.

For flag tests, $Q$ and any large $q$ are affine combinations of flag
nodes and bounded coordinates, with fixed real coefficients on the large
nodes and positive leading coefficients.  Convert every participating
primary basis node to a dependent log node.  Also convert any free node
in the resulting log formulas that will be used as a polynomial factor.
Saturate as in Section~\ref{sec:flags}.  Consequently every required power
of a primary node, and every finite polynomial in its log formula, is a
finite sum of shifts times bounded analytic factors.  Subsequent finite
flag refinements are permitted.

\begin{theorem}[Additive primitive packets]\label{thm:additive-packets}
Under these assumptions, the output maps of \eqref{eq:additive-A} and
\eqref{eq:additive-B} satisfy the packet conditions of
Definition~\ref{def:packet} and the differentiated primitive conditions of
Assumption~\ref{ass:calculus-primitive}.  They have no exceptional column
loss: every working column is safe.  Their only nonholomorphic defects
above the first band have simple costs $e^{\Re f_i}$ and, when $q$ is
large, $e^{\Re f_l}$, where $i,l$ are the leading indices of $Q,q$.
Each cost is used only while its log formula is retained.  The complete
terminal coefficient germs, on each lateral determination, are unchanged
by inserting unused nodes and by retesting a fixed finite construction
on a permitted refinement preserving the old active free nodes, with the
same lifted relations and limiting ordinary data.
In particular this primitive library is retestable for elimination
in the sense of Definition~\ref{def:elimination-retestable}, with the
fixed local lateral determinations.
\end{theorem}

We prove the theorem in several steps.  Up to and including band $i$ the
packet is the actual passage, with zero prefixes at earlier transitions.
The flag estimates give $\Re Q\gg\log|Q|$ and
$\Re(Q-q)\gg\log|Q|$; on the $i$ fringe, after increasing the side
dominance constant,
\begin{equation}\label{eq:additive-charge-realpart}
 \Re(Q-q)\ge c\Re Z_i.
\end{equation}
On the upper determination the endpoints lie in slightly enlarged first
quadrants and $|q/Q|$ is small; the lower determination is reflected.
For bounded $q$ use its ordinary neighborhood near the positive ray.

In the $A$ plane join $q$ horizontally to $q+S$, then straight to $Q-S$,
then horizontally to $Q$, where $S=C\log|Q|$ and $C$ is a sufficiently
large fixed number.  The real part increases.  On this polygon
$|A|\gtrsim|q|$, and on its middle segment $|A|$ bounds below a fixed
multiple of the convex combination of the endpoint moduli.  The implicit
formula for $w(A)$ gives
\[
 w(A)=A(1+O(|\delta|)),\qquad |\dd w/\dd A|\le C,
\]
because $(q/A)\log(A/q)$ is bounded there.  In particular
\begin{align}
 \int |X|\,|\dd w/w|&\le C\max_j|\chi_j|,
       \label{eq:additive-slow-integral}\\
 \sup(|E|+|H|)+\int(|E|+|H|)\,|\dd w|
       &\le C\max_s(|e_s|,|d_s|).\label{eq:additive-fast-integral}
\end{align}
Include $|\delta|$ in the first maximum.  For example, a positive slow
power integrates as $\int_0^1 t^{r-1}\dd t=1/r$; a negative power has
the corresponding inner-endpoint estimate.  On an endpoint horizontal
segment each active fast profile is an exponential in horizontal
distance times a fixed polynomial.  On the middle segment the margins
$S$ pay every fixed power and the length.  These observations prove the
two inequalities on all three segments, with constants depending only
on the fixed exponents and boxes.

The exact homogeneous propagator in (A) is
$e^{\sigma(A(w)-A(w'))}$; its modulus is at most one in the prescribed
direction.  The integral equations, \eqref{eq:additive-slow-integral},
\eqref{eq:additive-fast-integral}, and smallness keep solutions strictly
inside their state boxes.  The same estimates hold with parameter room.
Locally freeze the horizontal lengths and vary the endpoints.  Analytic
ODE dependence gives holomorphy; changing lengths by fixed factors gives
the same solution by a homotopy through the bounded paths.  This supplies
one continuation from the real analytic family, including across the
real ray.  No sector normalization has yet been introduced.

The sectorial constructions below normalize auxiliary solutions, while
the actual passage just constructed is fixed.  A simple instance of (A)
shows why their two endpoint contributions must be kept together.  Take
$\delta=0$, $g(X,v)=X=\chi q/w$, and the damped direction.  A particular
solution $V$ of $V'+V=\chi q/w^2$ has the formal expansion
\[
 \widehat V(w)=\chi q\sum_{j\ge0}(j+1)!w^{-j-2}.
\]
Two lateral particular solutions can differ by $Ke^{-w}$, whereas the
actual transfer is
\begin{equation}\label{eq:additive-Stokes-cancellation}
 V(p)+P\{v_{\rm in}-V(q)\}.
\end{equation}
Here $p=Q$, so changing $V$ by $Ke^{-w}$ changes the displayed value
by $Ke^{-p}-PKe^{-q}=0$.  The exact inner value $V(q)$ must therefore
remain in the coefficient of $P$.  Replacing it independently by its
formal series would discard the term that cancels the outer lateral
change.  The general construction retains inner endpoint values for
this reason.

\subsection{Lateral domains and the slow equations}

After the $i$ transition, $Q$ can acquire a large lifted log phase.  The
slow equations will use its lifted powers while $q$ remains in its
lateral quadrant, or near its bounded positive value.  Freeze a
representative of $|Q|$, take $B$ a sufficiently large fixed multiple of
that representative, and use the upper polygonal domain
\begin{equation}\label{eq:additive-polygon}
 \Omega_B^+=\{x+iy:-B<x<B,\ |y|<M(B-|x|),\
                  y>c_0|q|-\eps_0|x|\}.
\end{equation}
Here $0<c_0\ll\eps_0$, $M$ is fixed and large, and admissible limiting
anchor vertices are $-B$ and $B$.  Choose the constants so that a slightly
enlarged half-annulus
\[
 |q|/2\le |w|\le 2|Q|,\qquad
 -\eta_0\le\arg w\le\pi+\eta_0
\]
has distance at least $\eps|w|$ from the edges.  Increasing $B/|Q|$
provides any fixed additional outer room.  The lower domain is its
reflection.  Each has a fixed logarithm, $|w|\gtrsim|q|$, and all the
profiles stay in a small polydisk if their amplitudes were chosen small
enough, also on the larger domain.

Every point of \eqref{eq:additive-polygon} can be reached from either
vertex by a graph monotone in $x$ with $|\dd w|\le C|\dd x|$.
Indeed the upper and lower boundaries are Lipschitz graphs with bounded
slopes; their vertical intervals collapse at the vertices, and
interpolation between boundary graphs gives the claimed paths.  Nearby
graphs are homotopic through such graphs.  Volterra solutions with the
appropriate damping direction therefore continue to one analytic
function on the domain: locally vary the final path segment and use
uniqueness to compare overlapping continuations.  Starting at a vertex
can equivalently be defined by a limit from a full interior graph;
the field is nonsingular at that nonzero-radius point and analytic with
room beyond it.  The finite predecessors in a coefficient recursion
extend to the same vertices in this fashion.  Parameter changes with
anchors frozen give local holomorphic choices.

For (A) let $v_0$ solve
\begin{equation}\label{eq:additive-graph}
 v_0'-\sigma(1+Z)v_0=w^{-1}g(X,v_0)
\end{equation}
with zero data at the decaying vertex.  Define
$b=g_v(X,v_0)$ and $\mathfrak a=\sigma(1+Z)+b/w$.
The fiber change and its equations are
\begin{align}
 v=T(w,u)&=v_0(w)+u+\sum_{m\ge2}t_m(w)u^m,
       &u'&=\mathfrak a u,\label{eq:additive-fiber}\\
 t_m'+(m-1)\mathfrak a t_m
   &=w^{-1}[g(X,T)-g(X,v_0)-b(T-v_0)]_{u^m}.
       \label{eq:additive-fiber-recursion}
\end{align}
The $t_m$ have zero data at the opposite vertex.  This opposite choice is
forced by the sign of their diagonal.

For (B), let $m,n$ be the multi-indices of $E,H$, respectively, and put
\begin{equation}\label{eq:additive-charges}
 a=\sum_s a_sm_s,\qquad b_f=\sum_s b_sn_s,\qquad
 \Delta=b_f-a,\qquad \kappa=\beta\cdot m+\nu\cdot n.
\end{equation}
Thus $a,b_f$ are scalar total charges.  Seek
$z=u+t(w,E,H,u)$ with $t$ supported on $\Delta\ne0$ and
$u'=\mathcal C(w,E,H,u)$ with $\mathcal C$ supported on
$\Delta=0$, without a constant term.  Coefficient comparison gives
\begin{align}
 \mathcal C&=[G(X,E,H,u+t)]_{\Delta=0},\label{eq:additive-balanced}\\
 t_{mn}'+[\Delta(1+Z)+\kappa/w]t_{mn}
   &=[G(X,E,H,u+t)-\partial_ut\,\mathcal C]_{mn},
           \qquad\Delta\ne0.\label{eq:additive-imbalance}
\end{align}
There is no balanced term in $\partial_ut\,\mathcal C$: adding a
balanced charge to a nonzero imbalance preserves that imbalance.
These equations are triangular in total fast degree.  Each imbalance
has zero data at its decaying vertex, the left vertex when $\Delta>0$.
For a balanced monomial, the $w$-dependent fast exponentials cancel:
the corresponding term of the balanced equation is
\begin{equation}\label{eq:additive-balanced-P}
 P^j e^md^n q^{-\beta\cdot m}Q^{-\nu\cdot n}
          w^{\kappa}\mathcal C_{mn}(w,u),\qquad
 j=a=b_f\ge1.
\end{equation}
This is the drift that survives the removal of the imbalanced terms.

A fixed charge budget $K$ needs only coefficients modulo
$\min(a,b_f)>K$.  The retained region in
Figure~\ref{fig:retained-charges} is a finite union of fixed-side tails,
not a finite total-degree truncation.  The next lemma proves convergence
of each such tail on a common fast radius.  The actual remainder after
omitting the other charges is estimated later in
\eqref{eq:additive-charge-suppression}.

\begin{figure}[tb]
\centering
\begin{tikzpicture}[x=1.05cm,y=0.70cm,>=Stealth,font=\small]
  \fill[blue!9] (0,0) rectangle (2,5);
  \fill[blue!9] (2,0) rectangle (6,2);
  \draw[gray!60,dashed] (2,0) -- (2,5);
  \draw[gray!60,dashed] (0,2) -- (6,2);
  \draw[->] (0,0) -- (6.5,0) node[right] {$a$};
  \draw[->] (0,0) -- (0,5.5) node[above] {$b_f$};
  \node[below left] at (0,0) {$0$};
  \node[below] at (2,0) {$K$};
  \node[left] at (0,2) {$K$};
  \draw[blue!60!black,very thick] (0,0) -- (2,2);
  \fill[blue!60!black] (1,1) circle (1.6pt);
  \fill[blue!60!black] (2,2) circle (1.6pt);
  \node[rotate=34,anchor=south] at (1.05,0.96) {$a=b_f$};
  \draw[->,thick,blue!60!black] (1,2.6) -- (1,4.8);
  \node[align=left,anchor=west] at (1.15,4.3) {fixed $a\le K$};
  \draw[->,thick,blue!60!black] (2.7,1) -- (5.8,1);
  \node[anchor=south] at (4.2,1.05) {fixed $b_f\le K$};
  \node[align=center] at (4.2,3.3) {omitted\\$a>K$, $b_f>K$};
  \node[anchor=north] at (3,-0.55) {schematic charge region; not to scale};
\end{tikzpicture}
\caption{A finite request retains $\min(a,b_f)\le K$, a finite union
of convergent fixed-side tails; its balanced part is finite.
The excluded region has both charges larger than $K$, so its
monomials carry the corresponding power of the transfer charge on
the actual passage.  Not every schematic lattice position need occur.
No coefficient bound uniform in the fixed side, or convergence of
the unrestricted mixed series, is asserted.}
\label{fig:retained-charges}
\end{figure}


\begin{lemma}[Slow coefficients on a common fast radius]\label{lem:additive-slow}
The graph, fiber transformation in its convergent state-disk norm, and
$b$ in (A) are bounded on the lateral domains.  The transformation is
close to the identity and has a regular inverse on a smaller disk
covering the inputs.  In (B), for each fixed $m$ the full $H$ tail
converges on one fixed fast radius; the corresponding assertion holds
with the sides interchanged.  These radii cover the actual fast
amplitudes with fixed slack and do not decrease with $m,n$, or $K$.
State, profile, and spatial neighborhoods have common positive remaining
margins.  The finite balanced lists have the same estimates.

Each such fixed request has a formal series in $h=1/w$, whose coefficients
are analytic in independent profile and ordinary arguments on fixed
smaller polydisks and satisfy
\begin{equation}\label{eq:additive-Gevrey}
 \|c_j\|\le C^{j+1}j!.
\end{equation}
On padded interiors, for $1\le J\le c|w|$,
\begin{equation}\label{eq:additive-Gevrey-error}
 \left\|c(w)-\sum_{j<J}c_j(X)w^{-j}\right\|
       \le C^{J+1}J!|w|^{-J};
\end{equation}
a sufficiently small proportional optimal cutoff gives $O(e^{-c|w|})$.
For unscaled positive-profile coefficients
$\tau_j=\chi_jQ^{-r_j}$, so that $X_j=\tau_jw^{r_j}$,
the same estimates for $\partial_\tau^\alpha$ acquire the factor
$|w|^{r\cdot\alpha}$.  Fixed ordinary derivatives are allowed.
Neighboring anchor choices of size comparable to $|Q|$ differ by
$O(e^{-c|Q|})$ on the common working interior, including at its inner
radius, with fixed polynomial derivative losses allowed.
At $\tau=0$ the finite positive-profile jets have canonical
infinite-anchor determinations with these estimates.
\end{lemma}

\begin{proof}
We first construct the actual coefficients on the lateral domains.
For (B), this is an induction in the left degree: at each step a finite
list of opposite degrees precedes one convergent high opposite tail.
A fixed request involves only finitely many such blocks.
We then obtain formal Gevrey bounds and compare their optimal truncations
with those actual coefficients.  The common domain choices are made
before this second induction, so its slow order can grow without
exhausting the margins.  Weighted derivatives and anchor comparisons
complete the proof.  Throughout, a common fast radius does not assert
coefficient bounds uniform in the left degree or convergence when both
fast degrees grow.

\emph{1. Convergent lateral coefficients.}

For $\mathcal L=\partial_w+\Delta(1+Z)+\kappa/w$, its zero-data
inverse has kernel
\begin{equation}\label{eq:additive-kernel}
 \exp[-\Delta(w-w')-\Delta\gamma\log(w/w')]
                    (w'/w)^\kappa.
\end{equation}
On an admissible monotone graph, $|\gamma/w|\le C|\delta|$.
If $|\kappa/\Delta|$ is universally bounded, small $\delta$ and
large fixed $q_*$ give kernel decay $Ce^{-c|\Delta||x-x'|}$
and inverse norm $C/|\Delta|$.  For a fixed exceptional pair the last
factor in \eqref{eq:additive-kernel} is bounded by
$C(1+|x-x'|)^{C_\kappa}$: the log angle is bounded and both radii are
bounded below.  Its exponential moment is finite.  Thus exceptional
pairs only change their constants; they do not require increasing
$q_*$ as the requested charge order increases.  Fixed profile weights
commute through the kernel with the same polynomial ratio estimate.

The graph equation is a contraction on a small sup-norm ball, because
$g=O(X)$ and its state derivative is $O(X)$.
For the fiber coefficients use the norm $\sum_{m\ge2}R^m\|t_m\|$ on a
small state disk.  The inverse gains $C/(m-1)$, compensating the degree
factor in analytic composition.  More explicitly, expand $g$ in a disk
strictly inside its analytic state box, majorize its coefficients by
the corresponding absolutely convergent positive Taylor series, and
solve successive coefficient truncations in a small ball.  The
$1/w$ multiplier is at most $C/q_*$; terms containing $mb/w$ have
uniformly small norm after division by $m-1$.  The majorant map has
small Lipschitz constant.  Its bound is independent of the truncation,
so the increasing sums converge, and Cauchy on a smaller state disk
makes $T_u-1$ small.  The inverse follows by contraction.

The pure $E$ and pure $H$ transformations in (B) have no drift
interaction.  Their right sides have a fast factor.  A Taylor majorant
on small fixed fast disks and a larger state disk is therefore a
contraction with small state derivative.  Now fix $m\ne0$ and write
\[
 T_m(H)=\sum_n t_{mn}H^n,\qquad
 A_0(H)=G_z(X,0,H,u+T_0(H)).
\]
Here $A_0(0)=0$.  At the current left order, the unknown $T_m$ enters
$G$ only as the linear convolution $A_0T_m$; every nonlinear remainder
uses smaller componentwise left orders.  Also $\mathcal C_{0,*}=0$,
so $\partial_ut\,\mathcal C$ never differentiates the current unknown
$T_m$.  Current balanced coefficients multiply only
$\partial_uT_0$; they form a finite list with $b\cdot n=a\cdot m$ and
are determined in increasing opposite degree before the high tail.

Choose a fixed fast radius $\rho$ strictly inside the pure-side radius
and use the weighted absolute coefficient norm on a chosen lateral
domain $\Omega$ and state disk $D_R=\{|u-u_c|<R\}$,
\begin{equation}\label{eq:additive-tail-norm}
 \|F\|_{\rho,R}=\sum_n\rho^{|n|}
             \sup_{w\in\Omega,\,|u-u_c|<R}|F_n(w,u)|.
\end{equation}
It is a Banach algebra without any loss of $\rho$.  To make the uniformity
explicit, put $a_{\min}=\min_s a_s$, $b_{\min}=\min_s b_s$,
$\beta_{\max}=\max_s|\beta_s|$, and
$\nu_{\max}=\max_s|\nu_s|$.  For
$N=N(m)\ge\max(1,2a/b_{\min})$ and $|n|>N$,
\[
 \Delta_{mn}\ge\tfrac12 b_{\min}|n|,\qquad
 |\kappa_{mn}/\Delta_{mn}|
    \le\beta_{\max}/a_{\min}+2\nu_{\max}/b_{\min}.
\]
In particular the high-tail kernel constants and the initial choice of
$q_*$ are independent of $m$.
After solving the finite initial list, its high tail $U$ satisfies
\begin{equation}\label{eq:additive-tail-contraction}
 U=\mathcal L_m^{-1}\Pi_{>N}(A_0U+F_m),\qquad
 \|\mathcal L_m^{-1}\Pi_{>N}F\|_{\rho,R}
      \le\frac{C_0}{N}\|F\|_{\rho,R}.
\end{equation}
The forcing $F_m$ consists entirely of known lower-left tails and the
finite initial list.  Taking $N>2C_0\|A_0\|_{\rho,R}$ yields
\[
 \|U\|_{\rho,R}\le(2C_0/N)\|F_m\|_{\rho,R}.
\]
This proves convergence on the same $\rho$ at every left order.
Large mixed coefficients do not violate the field box: their Taylor
construction is coefficientwise in the left variables about the small
pure-side base, rather than an evaluation of the whole mixed series.

\emph{2. Domains shared by all slow orders.}
The preceding induction gives the actual coefficients on a common fast
radius.  To compare them with formal slow expansions, we must also reserve
state, profile, and spatial room at every finite dependency step.  We make
those choices now, independently of the slow truncation order.  Fix a
finite coefficient and derivative request and close its dependencies
under smaller left orders, opposite initial indices, and lower jets.
Choose positive limiting profile and state radii $x_\infty,R_\infty$
strictly inside the pure-side construction, and a spatial padding
$\zeta_\infty$ strictly inside the available geometric room.  Choose
$\eta_k=c_1/(k+1)^2$, $k\ge1$, with total sum smaller than each of the
three available margins, and define
\begin{equation}\label{eq:additive-reservoirs}
 x_k=x_\infty+\sum_{h>k}\eta_h,\qquad
 R_k=R_\infty+\sum_{h>k}\eta_h,\qquad
 \zeta_k=\zeta_\infty-\sum_{h>k}\eta_h.
\end{equation}
The fixed amplitudes can be chosen so that actual profiles lie in
$|X|<x_\infty/2$ on the whole larger lateral domain.  State coefficient
functions at degree $k$ are constructed on
$|u-u_c|<R_k+\eta_k/4$ and are estimated, with arbitrary fixed state
jets, on $D_{R_k}$.  Lower-left coefficients are available at $R_{k-1}$;
in particular, for a fixed derivative order $d$,
\begin{equation}\label{eq:additive-inherited-Cauchy}
 \|\partial_u^d F_{m'}\|_{\rho,R_k+\eta_k/4}
 \le d!\left(\frac{4}{3\eta_k}\right)^d
             \|F_{m'}\|_{\rho,R_{k-1}},\qquad |m'|<k.
\end{equation}
The same estimate applies to predecessor differences in a comparison.
The current unknown is never differentiated in $u$ by its defining
equation.  The reserved $\eta_k/4$ gives its requested output state
derivatives after it has been constructed.

There are finitely many left indices of degree $k$.  Take a common
threshold $N_k$ for their high opposite tails, first larger than all
$2a/b_{\min}$ and then large enough to absorb the pure multiplier in
both the actual and formal norms below.  List the finite opposite
degrees in increasing order, followed by the high tail.  Include the
finitely many lower-jet comparison steps required for the present
request.  Call the number of resulting steps $s_k$.  Each infinite
same-sign tail counts as one step, not as infinitely many indices.
For $0\le v\le s_k$, use the spatial domains
\begin{equation}\label{eq:additive-spatial-schedule}
 \Omega_{k,v}=\left\{w\in\Omega_B^\pm:
  \dist(w,\partial\Omega_B^\pm)>
   \left(\zeta_{k-1}+\frac{v\eta_k}{s_k}\right)|w|\right\}.
\end{equation}
These are estimates on restrictions of the same actual coefficient
functions, not newly prescribed boundary data.  A target in
$\Omega_{k,v}$ has, in either horizontal direction, a segment of length
$b_k|w|$ contained in $\Omega_{k,v-1}$, with all radii comparable to
$|w|$, where one may take
$b_k=\eta_k/[8s_k(1+\zeta_\infty)]$.
This follows from the $1$-Lipschitz property of distance and of modulus.
All orders end on the same positive remaining profile/state margins
and on a spatial padding less than $\zeta_\infty$.  Increasing a finite
jet request changes $s_k,b_k$ and the constants, without changing those
limiting margins or the fast radius $\rho$.

\emph{3. Formal Gevrey bounds.}
On the domains just chosen, write
\[
 \partial_w=h(D_X-h\partial_h),\qquad
 D_X=\sum_j r_jX_j\partial_{X_j}.
\]
Divide the recursions by their nonzero diagonal at $h=0$, a nonzero
integer charge times $1+Z$.  Work modulo $h^{J+1}$, $J\ge1$, put
$\theta=\varepsilon_k/(J+1)$, and, for $0\le j\le J$, set
\begin{align}
 x_{k,j}^{(J)}&=x_k+\frac{\eta_k}{4}
       +\frac{\eta_k}{4}\left(1-\frac{j}{J+1}\right),
             \label{eq:additive-formal-disks}\\
 \|F\|_{k,J}&=\sum_{j=0}^J\theta^j\sum_n\rho^{|n|}
     \sup_{\substack{|X|<x_{k,j}^{(J)}\\
                     |u-u_c|<R_k+\eta_k/4}}|F_{jn}(X,u)|.
             \label{eq:additive-formal-norm}
\end{align}
For a finite scalar coefficient omit the sum over $n$.  Every current
profile disk lies at least $\eta_k/2$ inside the predecessor target
disk $|X|<x_{k-1}$, and at least $\eta_k/4$ outside the current target
disk $|X|<x_k$.  Thus there is both inherited room and fixed residual
room.  In a product the target disk at degree $a+b$ lies in both input
disks, so this norm is an algebra.  The disk gap for a single slow
degree is $\eta_k/[4(J+1)]$.  Cauchy's inequality and the factor
$\theta$ give
\begin{equation}\label{eq:additive-truncated-operators}
 \|hD_X\|\le C\varepsilon_k/\eta_k,\qquad
 \|h^2\partial_h\|\le C\varepsilon_k,
 \qquad\|\Delta^{-1}h\kappa\|\le C\varepsilon_k
\end{equation}
on the high tail; the finite exceptional list has its own constants.

To spell out the induction inequality, let $M_{k,v}^{(J)}$ denote the
norm of the current formal coefficient block, and let $M_{\prec}^{(J)}$
be the maximum of the norms already constructed that enter its
forcing.  Fixed left-degree extraction of $G$ is a finite Taylor
calculation about the pure-side base.  Its terms have bounded analytic
multilinear coefficients and finitely many predecessor factors.
The drift contributes products with a lower-left state derivative,
bounded by \eqref{eq:additive-inherited-Cauchy}.  Hence there are fixed
$D_k,d_k$ and constants $C_{k,v}$, independent of $J$, for which that
forcing is bounded by
\[
 \Phi_{k,v}(M_{\prec}^{(J)})
 :=C_{k,v}(1+\eta_k^{-d_k})(1+M_{\prec}^{(J)})^{D_k}.
\]
The exponents are finite because the left Taylor degree and the jet
request are finite; no bound as $k$ increases is required.  Dividing
by the diagonal, and moving only its slow differential correction to
the right, gives for an exceptional coefficient
\begin{equation}\label{eq:additive-formal-induction}
 M_{k,v}^{(J)}\le
 \Phi_{k,v}(M_{\prec}^{(J)})+
 C_{k,v}\varepsilon_k(1+\eta_k^{-1})M_{k,v}^{(J)}.
\end{equation}
For the high tail it gives the same inequality with the last
coefficient replaced by
\begin{equation}\label{eq:additive-formal-tail-induction}
 \frac{C\,M_A}{N_k}
       +C_k\varepsilon_k(1+\eta_k^{-1}),
\end{equation}
where $M_A$ bounds the pure-side multiplier in its larger-disk formal
norm, uniformly in $J$.  Balanced blocks satisfy the forcing inequality
without a current unknown on the right.  The graph, fiber disk norm,
and pure-side blocks obey the corresponding analytic majorant
inequality with their already chosen small Lipschitz constant; their
slow differential correction has the same bound.

First choose $N_k$ to make $CM_A/N_k\le1/4$.  Induct through the finite
list, taking $\varepsilon_k$ smaller than every preceding admissible
choice and so small that all the differential coefficients in
\eqref{eq:additive-formal-induction}--\eqref{eq:additive-formal-tail-induction}
are at most $1/4$.  Absorption yields
$M_{k,v}^{(J)}\le2\Phi_{k,v}(M_{\prec}^{(J)})$.
Its finite recursion supplies a bound independent of $J$.
Shrinking $\varepsilon_k$ only decreases predecessor norms, so it
does not invalidate a prior bound.  This proves the simultaneous
formal induction on the specified disks.  At $J=j$ it implies
\[
 \|c_j\|\le C_0((j+1)/\varepsilon_k)^j
          \le C^{j+1}j!,
\]
using $j^j\le e^jj!$.  Fixed formal profile or state derivatives are
taken on the reserved target disks and change only their fixed
Cauchy constants.

\emph{4. Comparison with optimal truncations.}
The formal norm now bounds the substitution residual.  To turn this
into an estimate for the actual coefficient, we use a last horizontal
comparison segment of
length $\delta_m|w|$ in the appropriate damping direction, lying in
a padded interior, with radii comparable to $|w|$.  Substitute the
formal polynomial with $J$ a small multiple of $|w|$.  Test that
polynomial and all analytic operations at a dummy $h$ radius larger
than the actual $|1/w|$ by a fixed factor.  Formal cancellation below
degree $J$ and the just proved sum bound give an exponentially small
residual in the norm appropriate to its diagonal.  A diagonal or
$\kappa$ factor can contribute $1+|n|$ to this residual.  More precisely
introduce the forcing norm
\[
 \|F\|_{\rho,R,-1}
   =\sum_n\frac{\rho^{|n|}}{1+|n|}
       \sup_{\Omega\times D_R}|F_n|.
\]
The residual is $Ce^{-c|w|}$ in this norm, whereas the inverse maps
this forcing space boundedly into \eqref{eq:additive-tail-norm}.
Indeed apply the inverse coefficientwise before summing: on the high tail
\begin{equation}\label{eq:additive-index-cancellation}
 (1+|n|)/|\Delta_{mn}|\le C.
\end{equation}
Thus there is no hidden fast-radius loss in this step.

Here is the actual comparison inequality, including its starting term.
For a target of modulus $R$, take the segment from
\eqref{eq:additive-spatial-schedule}, oriented in the current damping
direction and parametrized by real distance $t\in[0,L]$, $L=b_kR$.
Let $E$ be actual coefficient minus its common formal polynomial.
The actual coefficients are uniformly bounded by the preceding
Volterra construction, and the formal sum norm bounds that polynomial
on the segment.  Thus
$\sum_n\rho^{|n|}\sup_{D_{R_k}}|E_n(0)|\le B_{k,v}$,
with no dependence on $R$ or $J$.  Normalize by $R^{r\cdot\alpha}$
for a fixed positive-profile jet; moduli on this segment are comparable
to $R$, so the same assertion holds with a different fixed constant.

Use the coefficient-sum weighted norms
\begin{equation}\label{eq:additive-weighted-segment}
 \|E\|_\omega=\sum_n\rho^{|n|}
       \sup_{\substack{0\le t\le L\\u\in D_{R_k}}}
                   e^{\omega t}|E_n(t,u)|,
 \quad
 \|F\|_{\omega,-1}=\sum_n\frac{\rho^{|n|}}{1+|n|}
       \sup_{\substack{0\le t\le L\\u\in D_{R_k}}}
                   e^{\omega t}|F_n(t,u)|.
\end{equation}
The sum is outside each coefficient supremum.  This convention is
essential for the inverse from the weak forcing norm to the strong
norm; no interchange with a single supremum of the weak sum is used.
The high-tail kernels are bounded by
$Ce^{-\lambda(1+|n|)(t-s)}$.  For $0<\omega<\lambda/2$, coefficientwise
integration therefore gives a bounded map from the second norm to the
first, and a bound $C/N_k$ from the strong forcing norm to the first.
Finite exceptional kernels have bounds $C_{k,v}e^{-\lambda_{k,v}(t-s)}$
and the analogous scalar assertion for $\omega<\lambda_{k,v}/2$.

At a current high tail the difference equation has the form
$\mathcal L_mE=A_0E+F$, where $F$ is the formal substitution residual
plus predecessor differences.  Precisely, if $P_J$ is the current
polynomial, $A_J,F_J$ are the polynomial predecessor evaluations, and
$\mathcal R_J=\mathcal L_mP_J-A_JP_J-F_J$, then
\[
 F=(A_0-A_J)P_J+(F_{\rm actual}-F_J)-\mathcal R_J.
\]
The finite Taylor operations defining
the forcing are Lipschitz on their bounded coefficient balls; their
Lipschitz constants are finite derivatives of $\Phi_{k,v}$.
Inherited state derivatives use
\eqref{eq:additive-inherited-Cauchy}.  Consequently, by the previous
comparison stages and the residual estimate, on this whole segment
\[
 \|F\|_{0,-1}\le C_{k,v}e^{-a_{k,v}R}
\]
for some $a_{k,v}>0$.  Predecessors may all use the current smaller
proportional cutoff: subtracting the finite Gevrey tails preserves
an exponential error.  In the exceptional list there is no current
multiplier; for a graph or pure side its place is taken by the small
linearized coefficient.  Variation of constants now gives
\begin{equation}\label{eq:additive-actual-induction}
 \|E\|_\omega\le C_{k,v}B_{k,v}
       +\tfrac12\|E\|_\omega
       +C_{k,v}e^{\omega L-a_{k,v}R}.
\end{equation}
The factor $1/2$ is ensured by the tail threshold, or the original
pure-side smallness, after the harmless weighted-kernel change.
Choose additionally $\omega b_k\le a_{k,v}/2$.  At the endpoint this
proves the explicit estimate
\begin{equation}\label{eq:additive-actual-endpoint}
 \|E(L)\|_{\rho,R_k}\le
 C_{k,v}\bigl(B_{k,v}e^{-\omega b_kR}+e^{-a_{k,v}R}\bigr).
\end{equation}
It has a positive exponential rate at every finite step.  Alternating
signs use the successive domains in
\eqref{eq:additive-spatial-schedule}; one infinite high tail uses one
segment, with its whole coefficient sum.  Thus the finite induction
closes without any fast-radius loss or exhaustion of the prescribed
state, profile, or spatial margins.  This proves the optimal error;
subtracting the omitted Gevrey terms proves
\eqref{eq:additive-Gevrey-error} at smaller orders.  Remaining bounded
radii for a fixed $J$ are absorbed in its constant.

\emph{5. Weighted profile derivatives.}
Differentiating at fixed anchors in $\tau$ gives polynomial weights
$w^{r\cdot\alpha}$.  The highest derivative of a graph or pure side
has the same small linearization, and the highest derivative of a high
tail has the same absorbable convolution.  The other terms use lower
derivatives.  Apply the unweighted resolvent first and then commute a
fixed polynomial weight through its decaying kernel; the resulting
constant is an exponential moment
$\int_0^\infty e^{-ct}(1+t)^{C_\alpha}\dd t<\infty$.
One need not demand that this weighted constant remain small.  More
explicitly, $A_0(0)=0$ means that each convolution step increases the
opposite degree.  In the iterated Volterra resolvent, a column starting
at degree $n_0$ is therefore bounded in weighted absolute coefficient
sum by
\[
 C e^{-\lambda(1+|n_0|)t}
       \sum_{v\ge0}\frac{(CM_A t)^v}{v!}
 \le C e^{-\lambda'(1+|n_0|)t},
\]
after the chosen high-tail threshold is increased once, if necessary.
For a pure side the same inequality follows from its original
smallness.  Along a monotone path, the ratio of either endpoint modulus
to the other is at most $1+Ct/q_*$.  Multiplying this resolved kernel
by a fixed power of that ratio gives a finite exponential moment and
still an inverse bound $C_\alpha/(1+|n_0|)$.  Thus the global weighted
jet bounds, and in particular the normalized starting bounds in
\eqref{eq:additive-actual-induction}, follow before any jet comparison.
Adding the finitely many lower-jet comparisons to the spatial schedule
then proves the differentiated errors with the same remaining domains.

\emph{6. Anchor comparison and limiting jets.}
Compare neighboring anchors using a common inner polygon with
horizontal extremes fixed large multiples of $|Q|$, leaving the
working target region $|\Re w|\le C_1|Q|$ strictly inside.  At each of
the finitely many comparison steps divide its allotted horizontal
margin by $s_k$ and choose a starting line on the appropriate side.
The monotone graphs used in constructing the lateral domain connect
that line to every target, with bounded slope and real distance at
least $b_{k,v}'|Q|$, where $b_{k,v}'>0$.  These paths can pass above the
inner obstacle; they need not be horizontal across it.  Graph and
pure-side differences have homogeneous absorbable equations.  At a
later step the difference forcing is bounded by the derivative of the
same finite polynomial majorant times predecessor differences.
The weighted argument \eqref{eq:additive-actual-induction}, now with
$R=|Q|$, therefore gives
\[
 E_{k,v}\le C_{k,v}|Q|^{p_{k,v}}
 \left(e^{-\omega b_{k,v}'|Q|}
                 +\max_{\prec}E_{\prec}\right).
\]
Here $p_{k,v}$ accounts for the fixed jet weights on the whole path;
without such jets it can be zero.  Finite induction proves
$E_{k,v}\le C_{k,v}'e^{-c_{k,v}|Q|}$ with $c_{k,v}>0$.
The starting distance is of outer-scale size also when the target is
near $q$, so the estimate holds uniformly down to the inner radius.
At $\tau=0$
positive-growth profiles disappear; their finite jets solve equations
with fixed polynomial weights.  Compare consecutive dyadic outer
anchors and sum their exponentially small differences on each fixed
interior.  The locally uniform limit gives the asserted analytic,
canonically normalized infinite-anchor jets and all their estimates.
\end{proof}

\subsection{Slow primitives and their normalization}

The linear coordinate in (A) needs a primitive $I$ satisfying
\begin{equation}\label{eq:additive-I}
 D I=b(w),\qquad D=w\partial_w.
\end{equation}
In (B), every balanced drift term in
\eqref{eq:additive-balanced-P} has positive degree in $P$.  A fixed
charge coefficient of its flow therefore uses only finitely many
iterations of the integral equation.  These give nested integrals of
entries $w^{d_e}c_e(w)$ in logarithmic time, with fixed rational $d_e$.
The factor changing $\dd w$ to $\dd w/w$ is included in $d_e$.
The functions $c_e$ are balanced slow coefficients or fixed derivatives
in an ordinary state argument, which is held independent at this step.
For each required word use a strictly upper triangular matrix having
its successive entries on the first upper diagonal.  A unit upper
triangular fundamental matrix $Y$ has entries satisfying
\begin{equation}\label{eq:additive-word}
 D Y_{\mathbf e}=w^{d_e}c_e(w)Y_{\rm tail},\qquad
 Y_\emptyset=1.
\end{equation}
Include all consecutive subwords.  Its transfer is $Y(p)Y(q)^{-1}$,
a polynomial in these finitely many entries.  This representation uses
no exponential bound on an unrestricted matrix flow over a long
logarithmic interval.

\begin{lemma}[Normalized slow primitives]\label{lem:additive-primitives}
There are particular formal primitives
\begin{equation}\label{eq:additive-formal-word}
 \widehat Y_{\mathbf e}
 =w^{d_{\mathbf e}}\sum_{j\ge0}w^{-j}B_{\mathbf e,j}(X,L),
 \qquad d_{\mathbf e}=\sum_{e\in\mathbf e}d_e,
 \qquad L=\log(w/q),
\end{equation}
where $B_{\mathbf e,j}$ is polynomial in $L$ of degree at most the
word length and its analytic coefficient norm is at most $C^{j+1}j!$.
The analogous assertion holds for $\widehat I$ with $d=0$ and length one.

Normalize the actual primitives by a sufficiently small proportional
optimal truncation at a positive $w_*$ comparable to $|Q|$, in the common
padded interior.  Then $I$ is bounded on the needed paths, and on a
matching outer ray the normalized outer values are approximated by
their formal optimal truncations with error $O(e^{-c|Q|})$.
The inner kernels
\begin{equation}\label{eq:additive-inner-kernels}
 I(q),\qquad q^{-d_{\mathbf e}}Y_{\mathbf e}(q)
\end{equation}
are bounded, have Gevrey expansions in $1/q$ obtained by setting $L=0$
in \eqref{eq:additive-formal-word}, and obey the differentiated estimates
of Lemma~\ref{lem:additive-slow}.  In particular a positive-profile jet
of order $\alpha$ is bounded after division by $q^{r\cdot\alpha}$.
Neighboring finite normalizations change these kernels by
$O(e^{-c|Q|})$, allowing fixed polynomial losses.

At zero positive profiles there are canonical infinite-anchor primitive
jets.  Their normalized inner values have the same Gevrey estimates.
For each fixed jet and each fixed $N$, finite and infinite normalized
jets differ by $O_N(|Q|^{-N})$.  This last assertion is power accuracy,
and does not claim exponential accuracy in $|Q|$.
\end{lemma}

\begin{proof}
\emph{1. Formal primitives and outer normalization.}
Multiply the already constructed formal series and argue by word
length.  The factorial inequality
$\sum_{k=0}^j k!(j-k)!\le Cj!$ preserves the asserted Gevrey order.
On a profile monomial $X^\mu$ the primitive operation is inversion of
\begin{equation}\label{eq:additive-resonance-operator}
 r\cdot\mu+d_{\mathbf e}-j+\partial_L.
\end{equation}
All its nonzero constant parts have absolute value at least $1/D_0$
for one fixed common rational denominator $D_0$.  On polynomials of
degree at most $s$, for $c\ne0$ its inverse is the finite expression
\[
 (c+\partial_L)^{-1}F
       =\sum_{k=0}^{s}(-1)^kc^{-k-1}\partial_L^kF.
\]
At $c=0$ integrate in $L$, choosing zero integration constant.  These
operators have bounds depending only on $s,D_0$.  Use absolutely
convergent profile coefficient norms on smaller disks to sum them.
Induction gives \eqref{eq:additive-formal-word} with degree at most
the word length.  In particular the formal choices have no undetermined
constants.

Freeze $w_*$ and the truncation order on each modulus choice.  Define
the actual triangular solution by those truncated values at $w_*$;
define $I$ in the same way.  For its leading coefficient,
$\widehat b_0=g_v(X,0)$ has no constant profile monomial.  A resonant
monomial at slow order zero must therefore involve both positive and
negative profile weights.  If the total positive weight is $\rho>0$,
its evaluation contains $(q/Q)^\rho$, and these positive weights have
a fixed positive rational minimum.  Thus the coefficient of $L$ at
order zero pays $\log(w_*/q)$.  Higher slow orders carry at least
$w_*^{-1}$, which pays this logarithm as well.  The normalized value of
$I$ at $w_*$ is bounded.  Since $b=O(X)$ on the actual graph, integrating
along radial and bounded-angle arc paths gives a bounded variation by
the nonzero profile powers, as in \eqref{eq:additive-slow-integral}.
This proves the boundedness needed for its later exponential.

If $p$ lies on the matching $Q$ ray, connect it to $w_*$ using paths with
all radii comparable to $|Q|$.  A common smaller optimal truncation
has exponentially small substitution residual, by
Lemma~\ref{lem:additive-slow}.  Triangular integration loses only fixed
powers and fixed-degree logarithms; these are absorbed by decreasing
the exponential constant.  This proves the outer endpoint assertion.
No actual evaluation of $I(p)$ at a multiply wound $Q$ is required.

\emph{2. Uniform propagation to the inner endpoint.}
The outer normalization alone does not bound the inner kernels when
$|Q|/|q|$ diverges.  We compare the actual and formal primitives along
the full inward path.  The triangular equations preserve their explicit
power weights, and a sufficiently high slow order makes those weights
integrable with constants independent of the endpoint ratio.

Fix the finite word and jet list, closed under
consecutive subwords and lower jets.  Put $q_0=|q|$, $R=w_*$,
$p_\alpha=r\cdot\alpha$, $\beta_{v,\alpha}=d_v+p_\alpha$, and
$\ell(t)=1+\log(t/q_0)$.  Fixed ordinary derivatives may be included
in the jet list with weight zero.  First compare on the radial part
$q_0\le t\le R$, and then on a bounded-angle arc to $q$.
For a word $v=e\,v'$ define
\[
 P_v^J(w)=w^{d_v}\sum_{j<J}B_{v,j}(X,L)w^{-j},\qquad
 c_e^{<L_0}(w)=\sum_{a<L_0}c_{e,a}(X)w^{-a}.
\]
Formal coefficient comparison gives the exact finite identity
\[
 DP_v^J=w^{d_v}\sum_{a+b<J}c_{e,a}B_{v',b}w^{-a-b}.
\]
Consequently the local substitution residual is exactly
\begin{equation}\label{eq:additive-primitive-residual-identity}
 \mathcal R_v^J:=w^{d_e}c_eP_{v'}^J-DP_v^J
 =w^{d_v}\sum_{b<J}w^{-b}B_{v',b}
                  (c_e-c_e^{<J-b}).
\end{equation}
Thus its estimate does not presuppose an estimate for an actual
primitive remainder.  On the fixed smaller profile disks the formal
coefficient bounds and Cauchy's inequality give
\[
 |\partial_\tau^\eta B_{v',b}|
 \le A_\eta C_0^{b+1}b!\,|w|^{p_\eta}\ell(|w|)^S,
\]
with one fixed $S$ covering all log degrees in the request.
Lemma~\ref{lem:additive-slow} gives, for $1\le L_0\le c|w|$,
\[
 |\partial_\tau^\eta(c_e-c_e^{<L_0})|
 \le A_\eta C_0^{L_0+1}L_0!\,|w|^{p_\eta-L_0}.
\]
The derivatives are at fixed $q,w_*$ and other independent arguments,
so they commute with $D$.  After Leibniz differentiation of
\eqref{eq:additive-primitive-residual-identity}, each term has power
$d_v-b+p_{\alpha-\eta}+p_\eta-(J-b)=\beta_{v,\alpha}-J$.
Using $\sum_{b<J}b!(J-b)!\le3J!$ and absorbing the finitely many jet
binomial factors proves
\begin{equation}\label{eq:additive-primitive-residual}
 |\partial_\tau^\alpha\mathcal R_v^J(w)|
 \le C_\alpha^{J+1}J!\,
       |w|^{\beta_{v,\alpha}-J}\ell(|w|)^S.
\end{equation}

The initial normalization tail has the very same weight.  Indeed
choose $K_R=\lfloor\varepsilon R\rfloor$ on larger parameter disks
such that $C_0K_R/R\le\theta<1$.  Fixed derivatives cost Cauchy
prefactors, while the coefficient-growth constant $C_0$ can be chosen
on those larger disks.  For $J\le cq_0\le K_R$, the geometric tail
estimate and $Y_v(R)=P_v^{K_R}(R)$ give
\begin{equation}\label{eq:additive-primitive-initial-tail}
 \left|\partial_\tau^\alpha(Y_v-P_v^J)(R)\right|
 \le\frac{A_\alpha}{1-\theta}C_0^{J+1}J!
        R^{\beta_{v,\alpha}-J}\ell(R)^S.
\end{equation}
The anchors and integer orders are frozen in these derivatives.

Write $E_{v,\alpha}^J=\partial_\tau^\alpha(Y_v-P_v^J)$.
Its differentiated equation is
\begin{equation}\label{eq:additive-primitive-error-system}
 D E_{v,\alpha}^J
 =w^{d_e}\sum_{\eta\le\alpha}\binom{\alpha}{\eta}
       (\partial_\tau^\eta c_e)E_{v',\alpha-\eta}^J
                   +\partial_\tau^\alpha\mathcal R_v^J.
\end{equation}
Since $|\partial_\tau^\eta c_e|\le C_\eta|w|^{p_\eta}$,
the power of each predecessor error on the right is exactly
$d_e+p_\eta+\beta_{v',\alpha-\eta}-J
=\beta_{v,\alpha}-J$.
Choose $J>1+\max_{v,\alpha}\beta_{v,\alpha}$, the maximum being over
the entire finite system.  For any such $\beta$ and $q_0\le t\le R$,
\begin{equation}\label{eq:additive-inward-integral}
 R^{\beta-J}\ell(R)^S+
 \int_t^R s^{\beta-J}\ell(s)^S\,\frac{\dd s}{s}
 \le C_S t^{\beta-J}\ell(t)^S.
\end{equation}
For the integral put $s=te^x$, use
$\ell(te^x)\le\ell(t)(1+x)$, and bound
$\int_0^\infty e^{-x}(1+x)^S\dd x$.
The boundary term follows from the bounded supremum of the same
function.  Both bounds are uniform in $t,R,q_0,J$; in particular the
initial tail cannot accumulate a factor depending on $R/q_0$.

For explicit triangular bookkeeping, let
\[
 X_{v,\alpha}^J=
 \sup_{q_0\le t\le R}
  \frac{|E_{v,\alpha}^J(t)|}
       {C_0^{J+1}J!t^{\beta_{v,\alpha}-J}\ell(t)^S}.
\]
Equations~\eqref{eq:additive-primitive-residual}--\eqref{eq:additive-inward-integral}
give
\begin{equation}\label{eq:additive-primitive-triangular-bound}
 X_{v,\alpha}^J\le B_{v,\alpha}
       +C_S\sum_{\eta\le\alpha}B_{e,\eta}
                         X_{v',\alpha-\eta}^J,
\end{equation}
after choosing the common $C_0$ to dominate the residual constants.
All $B$'s depend only on the fixed finite request and domain margins.
The right side contains strictly shorter words.  Finite triangular
induction bounds every $X_{v,\alpha}^J$ independently of $J,R,q_0$.
The bounded-angle arc has radius comparable to $q_0$ and bounded lifted
logarithm; integrating the same finite triangular system there changes
only these constants.  At $w=q$, divide by
$q^{d_v+p_\alpha}$ to obtain $C^{J+1}J!q_0^{-J}$.

To cover all requested orders $J\ge1$, choose a fixed integer
$B>1+\max_{v,\alpha}\beta_{v,\alpha}$ and first make the comparison
at $J+B$.  Decrease the target proportional constant so that
$J+B\le c_0q_0$ on the unbounded range.  The apparent extra factorial
is canceled by the extra radius power:
\begin{equation}\label{eq:additive-fixed-J-shift}
 \frac{(J+B)!}{J!}\,q_0^{-B}
       =\prod_{s=1}^{B}\frac{J+s}{q_0}\le c_0^B.
\end{equation}
Add back the fixed $B$ formal terms of degrees $J$ through $J+B-1$.
Their Gevrey bounds give the same $C^{J+1}J!q_0^{-J}$ estimate, with
an enlarged constant.  The remaining bounded range of $q_0$ is handled
by first integrating down to a fixed radius above this threshold and
then over the remaining compact radial interval and arc.  There are
only finitely many allowed $J$ on that range, and $q_0\ge q_*>1$;
all constants are finite and uniform there.  This changes no original
smallness condition or initial $q_*$.  The case $I$ is the length-one
version of the same calculation.

\emph{3. Finite and infinite normalizations.}
For two neighboring finite normalizations compare first near their
outer normalization points.  Their formal optimal values and the
nonprimitive coefficients differ by $O(e^{-c|Q|})$.
Propagate inward using the triangular difference equations.  Every
word and fixed jet loses at most a fixed power of $|Q|$, since $q$ is
bounded below, so the inner difference remains exponentially small.

At $\tau=0$ use the canonical nonprimitive infinite-anchor jets from
Lemma~\ref{lem:additive-slow}.  Subtract a differentiated formal
truncation of sufficiently high fixed order $J$ and solve the residual
equations by integration from positive-real infinity, then by an
interior arc when required.  Convergence follows from
\eqref{eq:additive-inward-integral}.  Indeed a fixed derivative
$\partial_\tau^\alpha$ at zero selects exactly the positive-profile
multi-degree $\alpha$; all remaining profile powers are negative.
Its $j$th formal term is therefore bounded by
$C_{\alpha,j}|w|^{d_{\mathbf e}+r\cdot\alpha-j}
(1+\log(|w|/|q|))^s$.  Taking $J$ beyond all these fixed power costs
makes the subsequently omitted terms tend to zero at infinity.
Increasing $J$ gives the same
solution: the difference solves the homogeneous differentiated
triangular recursion, and its first nonzero entry is constant with
limit zero at infinity; inductively every entry is zero.  The same
construction at $J\le c|q|$ gives its Gevrey estimate.  Locally uniform
convergence gives analytic parameter germs.  At $w_*$ compare this
infinite solution and the finite normalization at any sufficiently
large fixed $J$.  Their difference is smaller than any prescribed
power of $|Q|^{-1}$, after allowing the fixed polynomial/logarithmic
losses.  Propagate to $q$ by the triangular equations, using the
exponentially small nonprimitive differences on the common path.
Choosing $J$ still larger pays the finite propagation losses and proves
the claimed $O_N(|Q|^{-N})$ estimate.
\end{proof}

In particular every bounded inner kernel $K$ just constructed admits a
finite Taylor expansion in the positive profiles at their current values
$\tau_j=\chi_jQ^{-r_j}$.  For a prescribed finite total Taylor degree
$N$, the integral Taylor remainder and the weighted jet bounds give
\begin{equation}\label{eq:additive-positive-Taylor}
 K(\tau)=\sum_{|\alpha|<N}\frac{\chi^\alpha}{\alpha!}
       (q/Q)^{r\cdot\alpha}\,
       \big[q^{-r\cdot\alpha}\partial_\tau^\alpha K(0)\big]
       +O\big((|q|/|Q|)^{r_{\min}N}\big),
\end{equation}
where $r_{\min}$ is the smallest positive profile weight.  For normalized
word kernels $K$ already includes $q^{-d_{\mathbf e}}$.  The Taylor
segment remains in the original independent profile boxes.
Replacing each finite zero-profile jet in this finite sum by its
infinite-anchor jet has arbitrary additional $|Q|^{-N'}$ accuracy.
If there are no positive profiles, the Taylor step is vacuous: retain
the single zeroth jet, with the same finite-to-infinite replacement.

\begin{lemma}[Bounded analytic Gevrey operations]\label{lem:additive-gevrey-operations}
Finite sums and products, analytic operations at bounded arguments with
fixed room in their convergence boxes, and regular inverses preserve
the normalized Gevrey and jet conclusions above.  For two independent
slow variables the resulting formal coefficients have bounds
\begin{equation}\label{eq:additive-rectangular}
 \|c_{ab}\|\le C^{a+b+1}a!b!.
\end{equation}
Rectangular optimal evaluation has separate errors
$O(e^{-c|Q_0|})+O(e^{-c|q_0|})$; changing only one optimal cutoff has
the accuracy of its own variable.  These conclusions remain uniform
when other bounded analytic arguments vary independently on smaller
fixed disks, and under fixed derivatives on those disks.
\end{lemma}

\begin{proof}
For finite budgets $J,N$, test the truncated inputs at dummy radii
$\varepsilon/(1+J)$ and $\varepsilon/(1+N)$.  Their nonconstant
increments are small in absolute coefficient norm and their constant
values stay strictly inside the analytic boxes.  Absolute Taylor
majorants therefore bound the output independently of $J,N$.
Extracting its $(a,b)$ coefficient with $J=a,N=b$ gives
\eqref{eq:additive-rectangular}, using $a^a\le e^aa!$ and the analogous
inequality for $b$.  For a regular inverse the analytic implicit
theorem applies on the dummy disks, uniformly on the smaller argument
boxes; the resulting majorant is bounded there.  Actual inverses and
their approximants differ by the residual times a bounded inverse
Jacobian, with the same finite differentiated bounds.

At sufficiently small proportional optimal orders the dummy radii
exceed the actual values by a fixed factor.  Terms omitted in one
variable have geometrically small absolute sum in that variable,
uniformly in the other dummy disk.  This gives the two separate
exponential errors and proves the assertion when one cutoff alone is
changed.  At a replacement of a bounded kernel, ordinary Cauchy bounds
first pay its actual error; the formal calculation is then performed
on the replacement.  Fixed derivatives use smaller disks with room.
For infinite-anchor jets the ordinary finite chain rule commutes with
the locally uniform limiting construction, so these are the same
canonical jets as those computed after the operation.
\end{proof}

\subsection{Finite expansions in the transfer charge}

For (A), use $T$ and $I$ from one fixed lateral construction at both
endpoints.  Exact conjugacy gives the terminal value
\begin{equation}\label{eq:additive-A-transfer}
 T\left(w_{\rm out},
 P e^{I(w_{\rm out})-I(w_{\rm in})}
                  T^{-1}(w_{\rm in},v_{\rm in})\right).
\end{equation}
The exact linear coordinate stays in its disk along the passage because
of homogeneous damping and the bounded small variation of
$\int b\,\dd w/w$.  Expanding the final evaluation of $T$ at $u=0$ to any fixed
degree in $P$ has a remainder $C_K|P|^{K+1}$.  Keep the actual bounded
inner kernels at $q$; replace outer kernels at $p$ by their optimal
formal values, with exponential error in $|Q|$.

Inside $e^{\pm I(p)}$ the formal logarithm deserves care.  Write its
coefficient at slow order $j$ as $R_j(X)L(p)$.  The order-zero
resonant sum is $(q/Q)^\rho$ times an analytic function of small
rational-power ratios, for some fixed $\rho>0$, by the opposing-profile
argument in Lemma~\ref{lem:additive-primitives}.  Treat
$(q/Q)^\rho L(p)$ as a small analytic argument.  For $j\ge1$ factor
one $Q^{-1}L(p)$ and shift the remaining slow series by one fixed
index, retaining the bounded factors of $p/Q$.  This avoids an
unbounded logarithm as an analytic input or a shift of Gevrey index
depending on the requested order.  Lemma~\ref{lem:additive-gevrey-operations}
then applies to the exponential.

For (B), retain the coefficients of
\eqref{eq:additive-balanced}--\eqref{eq:additive-imbalance} with
$\min(a,b_f)\le K$ and denote the result by
$t^{[K]},\mathcal C^{[K]}$.  The two fixed-side parts of $t^{[K]}$
are genuine analytic sums on their opposite fast disks, as established
for the region in Figure~\ref{fig:retained-charges}.
Their pure parts $t_L,t_R$ are small, including their state derivatives.  Although
the other coefficient norms may depend on $K$, every mixed retained
part evaluated on the actual integration polygon is
$O(C_K|Q|^{C_K}|P|)$, with every requested fixed state derivative.

Here is the estimate used repeatedly.  For an analytic fast sum all of
whose monomials have $\min(a,b_f)\ge k$, select from each side a
sub-multi-index of charge in $[k,k+c_*]$, where $c_*$ is the largest
generator weight.  Since both $\Re(A-q)$ and $\Re(Q-A)$ are
nonnegative, the exponential product of those selected factors is at
most $|P|^k$.  Their fixed rational powers cost at most
$C_k|Q|^{C_k}$.  The remaining fast factors sum on the common disk,
using the fixed amplitude slack.  Thus
\begin{equation}\label{eq:additive-charge-suppression}
 |F(E(w),H(w))|
       \le C_k|Q|^{C_k}|P|^k\|F\|_{\rho}
\end{equation}
for that restricted support, and likewise on fixed smaller disks for
the needed derivatives.  No unrestricted double mixed series is used.

Invert $u\mapsto u+t^{[K]}(q,E,H,u)$ to impose the initial condition;
the pure-side derivative is small and the mixed part is exponentially
small by \eqref{eq:additive-charge-suppression}.  Solve
$u'=\mathcal C^{[K]}$.  Every term has at least one charge $P$, so
this is a small perturbation over the polygon despite fixed polynomial
length losses.  Then $z=u+t^{[K]}(w,E,H,u)$ has residual
\begin{equation}\label{eq:additive-B-residual}
 O(C_K|Q|^{C_K}|P|^{K+1}).
\end{equation}
To verify it with possibly large mixed coefficients, Taylor-expand $G$
about $u+t_L+t_R$ in the mixed part through degree $K$.  The actual
Taylor remainder already has \eqref{eq:additive-B-residual}.  The
finitely many products are analytic sums on smaller common radii.
Their coefficients of charge grade at most $K$ cancel by
\eqref{eq:additive-balanced}--\eqref{eq:additive-imbalance}.
Transport differentiation preserves the charges already present;
unused terms therefore come only from the analytic composition and
$\partial_ut\,\mathcal C$ products.  Bound them by
\eqref{eq:additive-charge-suppression}.  Termwise differentiation is
valid on the smaller disks.  Finally, $G_z=O(E,H)$ gives a bounded
integral of the state Lipschitz constant by
\eqref{eq:additive-fast-integral}.  The integral comparison with the
exact passage multiplies \eqref{eq:additive-B-residual} by at most
another fixed polynomial factor in $|Q|$.

The balanced terms have the exact form
\eqref{eq:additive-balanced-P}.  The endpoint fast arguments are exactly
\begin{align}
 E(q)&=e,& H_s(q)&=d_s(q/Q)^{\nu_s}P^{b_s},\label{eq:additive-end-fast-q}\\
 E_s(p)&=e_s(p/q)^{\beta_s}P^{a_s},&
 H_s(p)&=d_s(p/Q)^{\nu_s}.\label{eq:additive-end-fast-p}
\end{align}
Replace $P$ here and in \eqref{eq:additive-balanced-P} by a dummy
variable.  For $|P_{\rm dummy}|\le |Q|^{-C_K}$, with $C_K$ sufficiently
large, the initial inverse, balanced evolution, and final transformation
are analytic and bounded: all mixed terms and integrated balanced
terms are small, and the pure parts remain regular.  Cauchy's formula
therefore gives a Taylor expansion through degree $K$ with error
$C_K|Q|^{C_K}|P|^{K+1}$ at the actual exponentially small $P$.

Each coefficient in this expansion has a finite slow construction.
At degree zero the initial inverse is the pure-left inverse,
a bounded kernel $H_0(q,z_{\rm in})$ with regular endpoint estimates.
Higher inverse jets use fixed derivatives and the convergent
fixed-side tails.  Taylor--Picard coefficients of the balanced equation
are finite nested integrals of balanced coefficients and their state
derivatives at the constant $u=H_0$.  Splitting products of integrals
by the finitely many orderings of their integration variables writes
them as finite sums of the triangular transfer entries
$Y(p)Y(q)^{-1}$.  First evaluate these entries with $u$ independent;
then substitute $H_0$.  The final endpoint transformation uses the
same fixed-side coefficients and finite derivatives.  All power
prefactors in $q,Q,p$ are explicit and finite at each charge order.

\begin{proposition}[Finite charge expansion]\label{prop:additive-charge}
On any sufficiently tight $i$ fringe, through its buffered side, both
actual passages have an expansion to an arbitrary fixed charge cutoff,
with remainder
\begin{equation}\label{eq:additive-charge-remainder}
 C_K|Q|^{C_K}|P|^{K+1}+C_Ke^{-c_K|Q|}.
\end{equation}
Each coefficient is a finite expression in explicit power and
fixed-degree logarithmic prefactors, bounded normalized inner kernels
at $q$, and outer formal Gevrey series in a variable of size $Q^{-1}$.
Their analytic operations have bounded independent arguments with
fixed margins.  Outer coefficients may be evaluated at bounded inner
kernels such as $H_0$.  Every fixed derivative has the same description
with its own constants and polynomial losses.
\end{proposition}

\begin{proof}
The constructions above prove the expansion.  Formula
\eqref{eq:additive-endpoint-unit} writes $X(p)$, $p/Q$, and $1/p$ in
analytic units of the ratio and $(q/Q)\log(Q/q)$.  The formal word
coefficients have only fixed-degree log polynomials, which remain
outside analytic operations when unbounded.  The only exponential of
a primitive is controlled by the separate small arguments described
after \eqref{eq:additive-A-transfer}.  Thus the factorial estimates
apply to actual slow series indices; no estimate uniform in the charge
budget or in the finite explicit prefactors is being asserted.
For a fixed $M$, choose $K$ using
\eqref{eq:additive-charge-realpart}.  Its polynomial loss is negligible
relative to $\Re Z_i$.  Then tighten the fixed good contour until
$c_K|Q|$ dominates $M\Re Z_i$ throughout that fringe; proceeding to
the buffered side only improves this ratio.  The error is consequently
$O(e^{-M\Re Z_i})$.  The differentiated assertion follows from the
proved kernel derivatives, analytic operations on smaller disks, and
larger initial charge budgets.
\end{proof}

\subsection{From charge coefficients to packet trees}

Write the leading affine decompositions exactly as
\begin{equation}\label{eq:additive-leading-decomposition}
 Q=Q_0(1+\Theta_Q),\quad Q_0=A_0Z_i,\qquad
 q=q_0(1+\Theta_q),\quad q_0=A_1Z_l,
\end{equation}
where $A_0,A_1>0$ are fixed.  For bounded $q$ put $q_0=q$ and
$\Theta_q=0$, and regard $l$ as later than every transition.  Each
correction is a finite sum of strict positive shifts times bounded
analytic coefficients, with first index greater than $i$, or greater
than $l$, respectively.  This uses the stipulated conversion of
polynomial scale factors to log nodes.  Set
\begin{equation}\label{eq:additive-dispersal-variables}
 h_Q=Q_0^{-1},\qquad h_q=q_0^{-1}\quad(q\text{ large}),\qquad
 s_0=q_0/Q_0,\qquad M_0=\log(Q_0/q_0).
\end{equation}
The first index of $h_Q$ is $s=\sigma(i)$.  Let $r$ be the first index
of $M_0$, which has a positive leading coefficient.  Comparing the
exact log formulas gives
\begin{equation}\label{eq:additive-index-order}
 r\ge s,\qquad \sigma(l)\ge s\quad(q\text{ large});\qquad
 s\le l\ \Longrightarrow\ r=s.
\end{equation}
For bounded $q$, $r=s$.  Indeed no term of either log can occur before
$s$: a first term of $f_l$ before $s$ would contradict $f_i>f_l$.
If $s\le l$, the first index of $f_l$ is strictly later than $l$,
so it cannot cancel the leading term of $f_i$ at $s$.
Consequently
\begin{equation}\label{eq:additive-retirement-index}
 p_0=\min(l,r)\le s.
\end{equation}

All bounded formal outer arguments are analytic in finitely many
strict positive shifts with bounded coefficients.  Besides the $\Theta$
arguments, these include positive rational powers of $s_0$, the
components of $s_0^\rho M_0$ for fixed $\rho>0$, and those of
$h_QM_0$.  To see this last assertion, every unbounded term of $M_0$
is a node $Z_k$ with $k\ge r\ge s$.  Replacing this polynomial factor
by its log formula introduces only shifts with first index strictly
later than $k$.  Thus multiplication preserves the positive first
index $r$ of $s_0^\rho$, or $s$ of $h_Q$.  Logarithms of
$1+\Theta$ and of $p/Q$ are analytic near their regular values.
Explicit unbounded power and log prefactors outside analytic operations
are finite sums of shifts times analytic units in the same small
arguments.

In contrast to those analytic substitutions, the charge factor is
always extracted exactly:
\begin{equation}\label{eq:additive-exact-P-shift}
 P^n=\exp[-n(Q-q)].
\end{equation}
Its exponent is the original affine combination of nodes and bounded
coordinates, with leading term $nA_0Z_i$.  It is never evaluated by
Taylor-altering the exponent through $\Theta_Q$ or $\Theta_q$.

Retiring an anchor or cutoff means removing its numerical value from
subsequent coefficient formulas.  It does not identify its remainder
with zero.  Each retirement is performed only at a transition where
the stated remainder pays the requested exponential budget.  A
finite-anchor kernel, its optimal slow polynomial, and a fixed formal
coefficient are distinct objects; the exact lateral information has
already been retained in the charge construction, as illustrated by
\eqref{eq:additive-Stokes-cancellation}.

\begin{proposition}[Dispersal and retirement]\label{prop:additive-dispersal}
Every coefficient furnished by Proposition~\ref{prop:additive-charge}
admits, on each lateral determination, a packet tree using only retained
free inputs and ordinary inputs.  The numerical outer anchor for an
inner kernel is removed at $p_0$ in
\eqref{eq:additive-retirement-index}; the numerical optimal outer
variable is removed at $s$.  When $q$ is large, its optimal variable is
removed at $\sigma(l)$.  All transitions have arbitrary fixed
exponential accuracy on their full owning fringes.  The support is
iterated left-finite and all terminal leaves are analytic ordinary
germs.
\end{proposition}

The indices $l,r,s$ and $p_0=\min(l,r)\le s$ were defined in
\eqref{eq:additive-leading-decomposition}--\eqref{eq:additive-retirement-index}.
The retirement order is as follows.
\begin{center}
\small
\begin{tabular}{@{}>{\raggedright\arraybackslash}p{0.13\linewidth}>{\raggedright\arraybackslash}p{0.34\linewidth}>{\raggedright\arraybackslash}p{0.16\linewidth}>{\raggedright\arraybackslash}p{0.27\linewidth}@{}}
\toprule
Case & Finite outer anchor in the inner kernel & Outer optimal $h_Q$
 & Inner large-$q$ variable\\
\midrule
$l\le r$ & Removed at $l$ by optimal inner $q$ replacement
 & Removed at $s$ & Optimal $h_q$ removed at $\sigma(l)$\\[3pt]
$r<l$, $q$ large & Removed at $r=s$ by positive-profile Taylor
extraction and infinite-anchor jets
 & Removed at $s$ & Kernel replaced at $l$; optimal $h_q$ removed
at $\sigma(l)$\\[3pt]
$q$ bounded & Removed at $r=s$ by the same Taylor and jet construction
 & Removed at $s$ & None; ordinary infinite-anchor germs remain\\
\bottomrule
\end{tabular}
\end{center}
At coincident indices the formal replacement precedes extraction of
the current-index arguments; in particular, when $l=s$ the inner
replacement precedes extraction of $h_Q$.  In the last two rows the
finite-to-infinite primitive-jet comparison has arbitrary fixed power
accuracy in $Q$, as in Lemma~\ref{lem:additive-primitives}.  It is used
at $r=s$, where those powers pay the transition budget.

\begin{proof}
First set aside the finite explicit shift prefactors; a fixed shift
translates a packet tree.  We describe the remaining bounded
coefficient expression.  Keep its outer Gevrey series as an optimal
polynomial in $h_Q$ until index $s$ is extracted.  Rewrite analytic
unit factors into its coefficients, preserving $C^{a+1}a!$ bounds.
Keep the actual finite-anchor inner kernels through the band $p_0$,
using the respective upper or lower construction.

This prescription includes an independence requirement.  Before
$p_0$, put
\[
 Q_{\rm eval}=Q_0(1+\Theta_Q),\qquad
 \tau_j=\chi_jQ_{\rm eval}^{-r_j}\quad(r_j>0),
\]
using the ratio arguments currently retained or independently varied
in a Taylor calculation.  Freeze the anchor modulus and order using
$Q_0$, independently of these varied ratios.  Those ratios can vary
on fixed small complex disks: endpoint scales and profile sizes change
only by fixed factors, and the constructions have analytic room for
these changes.  Other artificial small arguments of the outer formal
coefficients need not vary the inner kernel.  Although the lifted
phase of $Q_0$ may wind, its needed logarithm is retained, and
$|Q_0|\gg|q|$.  All relevant powers have their fixed lifted branches.
At the first, $i$-matching evaluation these independent arguments take
their actual values.

At each transition, Taylor-expand every analytic small argument whose
first index is current to an arbitrary fixed budget.  Leave the other
small arguments independent; set an extracted one to zero in subsequent
coefficient evaluations.  This is permitted already before $p_0$.
Cauchy estimates on the original argument disks ensure that derivatives
created by such an earlier Taylor step retain the uniform kernel,
positive-profile, and Gevrey estimates needed at $p_0$.  In particular
an earlier Taylor expansion of a $Q$-ratio differentiates the bounded
holomorphic coefficient formula while $Q_0,q$ and the anchor convention
are held fixed.  It does not differentiate a varying anchor or
reimpose an identity between independent arguments.

There are two possible first retirements.

\emph{Case $l\le r$.}  At the $l$ fringe replace every bounded inner
kernel by its optimal $h_q$ expansion.  Its error $O(e^{-c|q|})$ pays
any fixed $l$-transition budget after the contour is tightened; split
past prefactors into their remaining shift tails before making this
comparison.  The formal coefficients use $X(q)$ and regular endpoint
arguments.  Positive profiles at $q$ carry powers of $s_0$ and hence
are among the allowed strict positive shifts.  Numerical finite-anchor
dependence has now disappeared.  If an outer coefficient was evaluated
at an inner kernel, use the rectangular series of
Lemma~\ref{lem:additive-gevrey-operations}, with coefficient bounds
$C^{a+b+1}a!b!$ in $h_Q^ah_q^b$.  Its separate optimal precisions
permit the two variables to be extracted at their different first
indices.  At index $l$ first perform this formal replacement and then
expand arguments whose first index is $l$; this may include $h_Q$,
whereas $h_q$ has first index $\sigma(l)>l$.

\emph{Case $r<l$, including bounded $q$.}  By
\eqref{eq:additive-index-order}, here $r=s$.  At that fringe use
\eqref{eq:additive-positive-Taylor} to Taylor-remove the positive
profiles and replace each retained zero-profile jet by its normalized
infinite-anchor jet.  Arbitrarily high fixed powers of $q/Q$ and
$Q^{-1}$ pay arbitrary depths at this transition: their logarithms
$M_0$ and $\log Q_0$ have positive first index $r$.  Keep the explicit
$s_0^{r\cdot\alpha}$ factors and analytic units outside the normalized
jets.  At this same index extract the required finite powers of $h_Q$.
One may do that finite $h_Q$ truncation first, uniformly in the bounded
inner arguments, then perform the inner Taylor and jet replacements.
Every resulting coefficient contains only bounded infinite-anchor
$q$-sector kernels and finite analytic operations.  It no longer uses
$Q_0$ numerically.  If $q$ is large, retain these kernels to the $l$
fringe and then replace them by their optimal $h_q$ series.  If $q$
is bounded, they are already ordinary analytic germs on the chosen
lateral determination.

For either case, a current-first optimal variable needs only a finite
degree at a fixed transition budget.  From its Gevrey bound and a
sufficiently small proportional cutoff,
\[
 \left|\sum_{j=J}^{J_{\rm opt}}c_jh^j\right|
          \le C_J|h|^J
\]
for fixed $J$.  Analytic Taylor tails satisfy the corresponding power
bound.  The flag side estimate makes every strict positive monomial
small throughout its available band, and at its owning fringe gives
a gain $e^{-c\Re Z_j}$.  Hence larger fixed $J$ gives every prescribed
transition precision.  Variables with later first index retain their
optimal evaluation uniformly on the independent parameter disks.
Split the finitely retained monomials at the current index, extract
their current powers, and leave later components as numerical
prefactors for the next band.  This recursively defines the packets.
The preceding argument applies separately to both angular signs.

This order of operations removes each numerical anchor or cutoff
before its controlling log formula becomes unavailable.  In the first
case the $Q$ anchor disappears at $l\le s$, and its remaining formal
optimal variable at $s$.  In the second case both disappear at $r=s$.
An inner optimal $q$ variable survives only until $\sigma(l)$.
The common-disk estimates in
Lemma~\ref{lem:additive-gevrey-operations} prove that each cutoff
change retains its own exponential precision, even in a rectangular
composition.  No error of merely power accuracy in $Q$ was used before
the transition $r=s$ where such powers pay the budget.

Finally the charge set at index $i$ is a nonnegative discrete lattice.
At a fixed charge there are finitely many prefactor translations,
followed by Taylor powers of a finite set of lexicographically strict
positive shifts.  At a fixed current-coordinate budget only finitely
many positive current-first factors can occur; zero current-first
factors are treated at the next coordinate.  After a prefix is fixed,
the remaining negative translations form a finite list.  Induction on
coordinates proves left-finiteness and finite collision multiplicity,
as in the support argument of Theorem~\ref{thm:calculus}.
The finite Taylor and jet operations at each later extraction preserve
this property.  At the terminal edge all numerical large arguments
have retired.  Formal coefficients and bounded-$q$ infinite-anchor
kernels are analytic ordinary germs, so every leaf is such a germ.
There is no assertion that the unrestricted two-sided fast series or
the full packet tree is simultaneously convergent.
\end{proof}

\subsection{Source precision and independent derivatives}

We finish by verifying the quantitative parts of the packet interface.
Choose anchors and proportional optimal orders deterministically on
fixed-factor bins of $A_0e^{\Re f_i}$ and, where needed,
$A_1e^{\Re f_l}$.  Locally freeze the representative and integer order.
Choose normalization systems for a fixed finite formula simultaneously
on larger parameter disks.  If an earlier coefficient has a different
small proportional convention, the overlapping comparisons permit its
fixed convention with only the already proved exponential error.
Smooth the local evaluations by partitions on overlapping retained
log-modulus bins.  All branches use the same lifted log relations.

Neighboring finite-anchor kernel evaluations differ by
$e^{-c|Q_0|}$ before their retirement.  Neighboring optimal polynomials
differ by $e^{-c|Q_0|}$ or $e^{-c|q_0|}$ while the corresponding
variable is active.  Since $|Z_k|=e^{\Re f_k}$, these are the simple
source costs
\begin{equation}\label{eq:additive-simple-costs}
 U_i=e^{\Re f_i},\qquad U_l=e^{\Re f_l}.
\end{equation}
They are used in a band $j$ only if the originating node $k<j$ still
has $\sigma(k)\ge j$.  Infinity jets and their bounded analytic
evaluations require no numerical anchor or cutoff and create no such
source.  The first band is the actual holomorphic family and is not
smoothed.

For completeness, the simple costs dominate all the losses even at
the buffered side.  Let $k<j$, $\sigma(k)\ge j$, write
$F=f_k(u)$, $L=f_j(u)$, and put $U=e^{\Re f_{k,j}}$ on band $j$.
We claim
\begin{equation}\label{eq:additive-source-dominance}
 U/|Z_j|\longrightarrow\infty
\end{equation}
uniformly to that side after increasing its finite dominance constant.
The following four cases use exactly the real, angular, and side
estimates of Section~\ref{sec:flags}.

If $F\le L^{3/2}$, the angular Taylor loss in the retained affine log
formula is
$O(F\operatorname{polylog}F/\Lambda_j^2)=o(1)$, since
$\Lambda_j\ge cL$.  Thus
$U/|Z_j|=e^{F-L+o(1)}\to\infty$.
If $L^{3/2}<F\le L^3$, then $\sigma(k)>j$, for a leading term at $j$
would make $F$ comparable to $e^L$.  The leading node in that log has
logarithm $\log F+O(1)=O(\log L)$, and angular speed bounded by a
fixed power of $\log L$.  Its phase on the $j$ band tends to zero.
The leading positive term consequently retains a fixed fraction of
its real size, and $\Re f_{k,j}\ge cF\gg L$.
If $F>L^3$ and $\sigma(k)>j$, the later-angle gain and side comparison
give
\[
 \Re f_{k,j}\ge cF/\Lambda_j
       \ge cF/(L\operatorname{polylog}L)\gg L.
\]
Indeed the phase deficit of the leading later node is at least a
constant times $1/\Lambda_j$, by the diverging gap derivative;
its later terms are negligible relative to that real leading term.
Finally if $\sigma(k)=j$, the exact log begins with a fixed positive
multiple of $Z_j$.  Increasing the side dominance constant absorbs
its finitely many later terms and gives
\[
 \Re f_{k,j}\ge c\Re Z_j
       \ge c e^L/L^{K+1}\gg L.
\]
These cases prove \eqref{eq:additive-source-dominance}.

After explicit prefactors have been set aside, the raw bounded pieces
have bounded size on every working band.  For example a resonant outer
logarithm is multiplied by its specified small positive shift; an
unbounded explicit log polynomial was kept outside analytic
operations and converted to finite shift prefactors.  A fixed remaining
shift uses indices at least $j$, so its log modulus is at most
$C\Re Z_j$ by side dominance.  Thus raw growth is
\begin{equation}\label{eq:additive-raw-growth}
 \log^+|S^j_\alpha|\le C_\alpha\Re Z_j
\end{equation}
throughout the band.  There are no sine poles or other exceptional
column factors.  Every column, and every needed entering collar, has
the safe improvement.  By \eqref{eq:additive-source-dominance}, all
these losses have log-size $o(U)$ at a source; products, interpolation,
and any fixed later normalization preserve $e^{-c'U}$ precision.

We verify derivatives in the independent retained free and ordinary
variables, rather than only along a path.  On band $j$, all numerical
node values of index at least $j$ and their finite triangular
sensitivities are bounded by $e^{C f_j(u)}$.  Any earlier retained log
is used through its exact affine/background expression in these
nodes.  With anchors and orders fixed, local Cauchy neighborhoods of
radius
\begin{equation}\label{eq:additive-Cauchy-radius}
 e^{-C f_j(u)}
\end{equation}
therefore preserve the small-argument disks, relative-modulus slack,
and angular room; enlarge $C$ for a fixed requested derivative order.
One first uses a slightly wider polynomial buffer and smaller
dominance constant, and then restricts to the working band.  Cauchy
differentiation loses at most $e^{C_Df_j(u)}$.  Since the side buffer
gives $\Re Z_j\gg f_j(u)^n$ for every fixed $n$, this loss is paid
by taking larger transition budgets.  At a source it has log-size
$o(U)$ by \eqref{eq:additive-source-dominance}.  Derivatives of smooth
partitions and of their costs have the same property: fixed powers of
$U$ and finite triangular sensitivity factors have logarithms $o(U)$.
Thus fixed derivatives of the antiholomorphic defects retain their
allowed source precision.  The difference
$\partial_{\Im z}-i\partial_{\Re z}$ vanishes for frozen holomorphic
choices and equals just these smoothing defects for the chosen packets.

The statements are stable under the controlled bounded-input tests
in Assumption~\ref{ass:calculus-primitive}.  To be explicit, use the
actual real controlled backgrounds for the real comparisons and their
common Taylor jets at the complex test point.  The triangular
sensitivities of a retained log have only a fixed power loss relative
to its leading size on the angular scale.  Taking a sufficiently high
fixed shared jet order makes the induced error smaller than the
available angular and polynomial side buffer.  In the small-$F$ case
of \eqref{eq:additive-source-dominance} take enough jets to leave an
absolute $o(1)$ log error; in the other three cases a sufficiently
small relative error preserves the displayed lower bound.  These
pointwise tests retain the local Cauchy neighborhoods with slightly
more slack.  The derivative, source, and safe-column assertions hence
hold simultaneously for any fixed finite collection of labels and
derivatives.  Starts, boxes, and constants may depend on that finite
request; every fixed further request is eventually available on the
same path returning to the ordinary germ.  No uniformity over all
budgets at one starting point is used.

Finite Taylor agreement locates the arguments in the common domains and
preserves the rate comparisons; it gives no source estimate for an
arbitrary smooth test.  For holomorphic inputs, or for a sheet already
having differentiated source bounds, apply
Lemma~\ref{lem:calculus-source-pullback}.  The primitive's intrinsic
defects have the simple costs in \eqref{eq:additive-simple-costs};
Equation~\eqref{eq:additive-source-dominance} pays their fixed losses.
Inherited sheet defects retain their own permitted costs, which may
include a cosine factor.  The general source comparisons accompanying
\eqref{eq:calculus-source-fringe} pay the same Cauchy and triangular
losses of logarithmic size $O(f_j)$, as well as the fixed raw and
normalization losses of size $O(\Re Z_j)$.  The chain-rule factors
therefore have logarithmic size $o(U)$ at the corresponding cutoff
sources in both cases.  The composed primitive retains the required
differentiated source bounds.

At a newly chosen level-$p$ good collar, a fixed current-scale
normalization has logarithmic loss at most
$C\Re Z_p\le C'\eps V_p$; fixed later-scale and derivative losses are
$o(V_p)$.  The join part of Lemma~\ref{lem:calculus-source-pullback}
therefore applies.  Its finite transition budget is chosen before the
contour, and errors $e^{-cV_p}$ with contour-independent $c>0$ allow the
stated small-$\eps$ choice.  This verifies the join precision separately
from the $o(U)$ estimates for active cutoff sources.

\subsection{Uniformity and canonical terminal germs}

We have obtained growth, transition, and differentiated source bounds.
We now track their dependence on the controlled input family, then show
that the terminal coefficients do not depend on the numerical anchors
and cutoffs used along the way.  Together with extension to preceding
contours, these properties permit the retests used in ordinary elimination.

\begin{lemma}[Uniformity over a fixed controlled family]
\label{lem:additive-controlled-family}
The additive constructions satisfy the common-constant requirement of
Definition~\ref{def:calculus-controlled-family}.  In particular, fix
one finite set of labels, derivatives, Taylor orders, and transition
budgets, a compact ordinary argument box with common positive analytic
slack, and common bounds for the finite real jets, Taylor errors,
angular rooms, and flag inequalities used by that request.  Then the
additive growth, transition, source, derivative, and safe-column
estimates have common constants and common numerical rate and gap
thresholds throughout every member's full forward domain where those
inequalities and argument ranges hold.  These constants and thresholds
are independent of the member's return time to the terminal germ.
\end{lemma}

\begin{proof}
First fix the field sup norms on the common larger complex boxes,
their distances to the working boxes, the amplitude bound, $q_*$,
the profile exponents and charge lattice, the polygon slope bounds,
and the finite requested coefficient and jet lists.  Include the
positive lower bounds for $|1+Z|$ and the regular inverse Jacobians
used in that finite construction.  These are finite quantitative
data on the stipulated compact boxes.  Every constant in the
$w$-integral equations is a function of these data: the parameters of
the ordinary test are frozen at its current evaluation while the
$w$-equation is solved.  Its earlier or later values do not enter the
Volterra equation.

In particular the reservoirs \eqref{eq:additive-reservoirs}, opposite
thresholds $N_k$, finite counts $s_k$, and polynomial majorants
$\Phi_{k,v}$ can all be chosen from these fixed data, before choosing
an ordinary test or its anchor.  Equations
\eqref{eq:additive-formal-induction}--\eqref{eq:additive-formal-tail-induction}
then give common $\varepsilon_k$ and formal bounds; Equations
\eqref{eq:additive-actual-induction}--\eqref{eq:additive-actual-endpoint}
give common positive comparison exponents.  The weighted resolvent
uses only the same coefficient norms and fixed polynomial jet powers.
For primitives, $\beta_{v,\alpha}$, $S$, the constants in
\eqref{eq:additive-primitive-triangular-bound}, and the fixed shift
$B$ in \eqref{eq:additive-fixed-J-shift} are likewise fixed by that
finite list.  Thus their remainder bounds and all analytic operations
have constants uniform over the independent argument boxes.  The
large-$Q$ and large-ratio qualifications used in this asymptotic model
are common numerical thresholds.  The bounded-$q$ alternative uses
its stipulated ordinary neighborhood; the compact radial intervals
used for small orders in the fixed-$J$ argument have the uniform
continuation bounds already proved.

Only passage from those independent arguments to a complex flag band
uses the ordinary test's real jets.  We give the quantitative check.
Suppose its background Taylor discrepancy satisfies
$|\Delta b|\le H_N|v|^{N+1}$ with the same $H_N$ for the family,
and $|v|\le C/\Lambda_j$ in band $j$.  The triangular sensitivity
bounds of Section~\ref{sec:flags}, on the common enlarged angular
domain, imply for retained $l\ge j$
\begin{equation}\label{eq:additive-family-later-log}
 |\Delta f_l|\le C H_N\Lambda_j^{-N}.
\end{equation}
For an earlier retained log $f_k$ with $\sigma(k)\ge j$, use its exact
affine formula in $Z_l$, $l\ge j$.  The sum of its term moduli is
bounded by a fixed multiple of $F=f_k(u)$, using its positive leading
term and the common gap thresholds.  Therefore
\begin{equation}\label{eq:additive-family-earlier-log}
 |\Delta f_k|\le C H_N F\Lambda_j^{-N}.
\end{equation}
Both estimates persist on the segment between the reference and test
backgrounds after the common angular room is narrowed, so their use
does not assume the desired test-domain conclusion.

Write $L=f_j(u)$.  In the small-log case $F\le L^{3/2}$,
\eqref{eq:additive-family-earlier-log} is an absolute $o(1)$ bound
uniformly in the family when $N>2$, since $\Lambda_j\ge cL$.
When $\sigma(k)>j$, compare it with the available lower bound
$cF/\Lambda_j$: its relative size is at most
$CH_N\Lambda_j^{1-N}$.  When $\sigma(k)=j$, compare instead with
the side lower bound $cF/L^{K+1}$: its relative size is at most
$CH_NL^{K+1}\Lambda_j^{-N}$.  Choosing the one fixed jet order
$N>K+2$ makes both ratios tend to zero with a common modulus.
Equation~\eqref{eq:additive-family-later-log} similarly preserves the
later-node angle slack and the polynomial side buffer.  This proves
the controlled-test versions of the domain conditions and all four
cases of \eqref{eq:additive-source-dominance} using only common jet
and flag constants.

The Cauchy neighborhoods \eqref{eq:additive-Cauchy-radius} can now use
one constant $C$ for the finite request.  The local bin partitions can
be taken from one fixed smooth partition in the retained log moduli;
their finite derivative bounds are numerical constants.  Multiplying
by the common triangular sensitivity bounds shows that all partition
and Cauchy losses have uniformly negligible logarithm relative to the
active simple costs.  There are only finitely many labels and costs in
this request, so their common minimum positive precision constant is
still positive.  Every column is safe, with the already uniform bound
\eqref{eq:additive-raw-growth}.

All arguments are pointwise quantitative implications of the finite
box, jet, rate, gap, and slack inequalities.  They apply at every
point of every full forward domain on which those common inequalities
hold, and not merely on a window selected separately for each test.
No step uses how soon a test returns to the germ, or makes a compactness
argument over such return times.  A later finite request can require a
smaller ordinary box and different controls; its constants are then
chosen from those new data, as permitted by
Definition~\ref{def:calculus-controlled-family}.
\end{proof}

Each fixed label, with its fixed derivative request, extends inward to
every sufficiently vertical preceding good contour.  Indeed its local
restrictions are small-argument inequalities and the bounded-angle
condition on its next still numerical endpoint.  The flag angular
comparison puts that endpoint in the slightly enlarged lateral quadrant
throughout this inward range; the remaining optimal expressions use
lifted logs with no further angle restriction.  All argument inequalities
have the fixed spare room used above.  The threshold for the preceding
contour may depend on the label and derivative order, as allowed by the
primitive contract.

It remains to record why the terminal germs are invariant in the
retesting used by ordinary elimination.  At a fixed charge, all formal
slow coefficients are fixed by the algebraic recursions with nonzero
diagonals; the formal primitive constants were fixed by
\eqref{eq:additive-resonance-operator}.  Charge coefficients of two
different finite budgets agree up to their common degree, because the
same triangular coefficient equations, regular initial inverse, and
finite Taylor--Picard recursion determine them.  Where a bounded inner
kernel survives, its zero-positive-profile jets are the canonical
infinite-anchor solutions fixed by decay of the subtracted remainder
at infinity.  Ordinary analytic compositions and regular inverses
commute with these finite coefficient and jet recursions.

Inserting an unused flag node therefore inserts an identity transition
with exponent zero.  A permitted refinement preserving old active free
nodes pulls back the monomial description, these same fixed recursions,
and the same ordinary germ functions;
only the grouping of intermediate coefficients changes.  Tiny free
perturbation sequences that preserve the affine log ties, regime gaps,
limiting ordinary data, and chosen lateral determinations consequently
give the same complete terminal coefficient arrays after this
identification.  Apply Theorem~\ref{thm:calculus} successively for any
fixed finite ordinary-chart recipe, always keeping its extracted
current symbol at the current inner background.  The terminal
operations are the same analytic operations and regular limiting root
equations, so the terminal germs remain the same.  This is a local
statement for the fixed lifted determinations; it claims no global
monodromy independence.  Nor does it identify a flat error with zero:
the charge construction, as \eqref{eq:additive-Stokes-cancellation}
illustrates, retains that information before retirement.

The full-fringe transition estimates of
Propositions~\ref{prop:additive-charge} and
\ref{prop:additive-dispersal}, the source and derivative estimates just
proved, and actual first-band holomorphy establish every assertion of
Theorem~\ref{thm:additive-packets}.
